\documentclass[a4paper,12pt]{article}

% Packages
\usepackage{float}
\usepackage{graphicx}   % for including images
\usepackage{setspace}   % for line spacing
\usepackage{geometry}   % for page margins
\usepackage{titling}    % for title customization
\usepackage{lipsum}     % for dummy text (can be removed)
\usepackage{newtxtext,newtxmath} % Times New Roman
\usepackage{xcolor}
\usepackage{algorithm}
\usepackage{algorithmic}
\usepackage{float}
\usepackage[normalem]{ulem}
\usepackage{calc} 
\usepackage{graphicx} % for including images
\usepackage{caption}  % optional, for better captions
\usepackage[style=ieee,backend=biber]{biblatex}
\addbibresource{references.bib}
\usepackage{listings}
\usepackage{xcolor}

\lstset{%
  language=Matlab,                % set language to MATLAB
  basicstyle=\ttfamily\footnotesize, % code font
  keywordstyle=\color{blue},      % keywords in blue
  stringstyle=\color{red},        % strings in red
  commentstyle=\color{green!50!black}, % comments in green
  numbers=left,                   % line numbers on the left
  numberstyle=\tiny\color{gray},  % style of numbers
  frame=single,                   % box around the code
  breaklines=true,                % wrap long lines
  tabsize=4                       % tab size
}

% Page setup
\geometry{top=1in, bottom=1in, left=1in, right=1in}
\begin{document}

% Title Section
\begin{center}
    \large \textbf{EXPERIMENT NO. 03-- Simulation and Performance Analysis of a Poisson Arrival–Poisson Service Queue with Finite Buffer}
    \rule{\linewidth}{1pt}
\end{center}

\vspace{3cm} % space before logo
\begin{center}
    \includegraphics[width=0.3\textwidth]{FrontPage/kgplogo.png} 
\end{center}
\vspace{2cm}
\begin{center}
    \textbf{\textit{Name: Dibyanshu Shekhar Dey}} \\
    \textbf{\textit{Roll Number: 25EC64R23}} \\
    \vspace{1cm}
    
\end{center}

\vfill % push bottom text to bottom of page

\begin{center}
    \rule{13cm}{1pt}\\
    \textbf{Design and Simulation Lab: Communication Networks} \\
    Department of Electronics \& Electrical Communication Engineering \\
    Indian Institute of Technology Kharagpur
    \rule{13cm}{1pt}
\end{center}
\rule{\textwidth}{2pt}

\newpage
%Start of the experiment
\noindent
\textbf{AIM:} To simulate a communication system with Poisson packet arrivals and Poisson service times by implementing a finite buffer and evaluate the system performance by evaluating the packet drop probability. \\
\noindent
\rule{\textwidth}{2pt}
\noindent
\textbf{INTRODUCTION: }\\
\noindent
In networks, queueing systems are used to analyse the network performance. In networks, packets arrive randomly, they are of random length. The main focus therefore shifts to delay related performance of the packets i.e- how long the packets would wait for and how much longer the buffer will be for the packets. This experiment tries to formulate a simulation so that the randomness can be captured for a network.\\

\begin{figure}[H]
    \centering
    \includegraphics[width=0.5\textwidth]{Introduction/buffer.drawio.png}
    \caption{Buffer and Serving Process}
    \label{fig:placeholder}
\end{figure}

\noindent
\textbf{THEORY: }\\
\noindent
The arrival process is modeled with Poisson process, which is a random process. In a Poisson process, the arrivals occur at an average rate of $\lambda$ packets per second \cite{ModianoQueueingMIT}. The probability of observing $n$ arrivals in a time period $T$ is given by \\
\[
            P(N(T) = n) = \frac{(\lambda T)^n e^{-\lambda T}}{n!}
\]

\noindent
The main idea in this experiment is to model the arrival process and service process based on Poisson process,. To do so, we use exponential random variable and the concept of inter arrival time. Inter arrival time denotes the time between two successive arrivals \cite{ModianoQueueingMIT}. Following is the CDF of inter arrival time: \\

\begin{align*}
    P(\text{IA} \le t) 
    &= 1 - P(\text{IA} > t) \\
    &= 1 - P(\text{0 arrivals in time } t) \\
    &= 1 - e^{-\lambda t}
\end{align*}
\noindent
The distribution is known as the \textbf{exponential distribution}. \\

\noindent
So. inter-arrival CDF:
        \[
            F_{\text{IA}}(t) = 1 - e^{-\lambda t}
        \]

\noindent
And inter-arrival PDF:
        \[
            f_{\text{IA}}(t) = \frac{d}{dt}F_{\text{IA}}(t) = \lambda e^{-\lambda t}
        \]

\begin{figure}[H]
    \centering
    \includegraphics[width=0.5\textwidth]{Theory/exponentialDistribution.png}
    \caption{Inter arrival Process Example (Arrival of N people in between $t$ time interval) \cite{ChanECE302ExpRV}}
    \label{fig:placeholder}
\end{figure}

\noindent
\textbf{Derivation of the Inter-arrival Time Distribution}
\noindent
Let $T$ denote the inter-arrival time between two consecutive events in a Poisson process with rate $\lambda$. Let $t$ be the time period under observation \cite{ChanECE302ExpRV}.

\begin{align*}
\mathbb{P}(T > t)
&\stackrel{(a)}{=} \mathbb{P}(\text{inter-arrival time} > t) \\
&\stackrel{(b)}{=} \mathbb{P}(\text{no arrivals in time } t) \\
&\stackrel{(c)}{=} \mathbb{P}(N(t) = 0) \\
&= \frac{(\lambda t)^0}{0!} e^{-\lambda t} \\
&= e^{-\lambda t}.
\end{align*}
\noindent
Therefore, the cumulative distribution function (CDF) of the inter-arrival time is
\[
F_{\text{IA}}(t) = \mathbb{P}(T \le t) = 1 - e^{-\lambda t}, \quad t \ge 0.
\]
\noindent
Differentiating the CDF yields the probability density function (PDF):
\[
f_{\text{IA}}(t) = \frac{d}{dt}F_{\text{IA}}(t) = \lambda e^{-\lambda t}, \quad t \ge 0.
\]
\noindent
Hence, the inter-arrival time $T = {\text{IA}}$ follows an exponential distribution with parameter $\lambda$.\\

\noindent
\textbf{RESULTS:}
\begin{verbatim}
Simulation Time           : 10.00
Arrival rate (lambda)     : 10.00
Service rate (mu)         : 13.00
Buffer size (K)           : 5
----------------------------------
Total Arrivals            : 96
Total Serviced            : 83
Total Dropped             : 8
Packet Drop Probability   : 0.083333
----------------------------------
Storing Interaarival Times
Inter-arrival[1] = 0.143
Inter-arrival[2] = 0.143
Inter-arrival[3] = 0.119
Inter-arrival[4] = 0.046
Inter-arrival[5] = 0.123
Inter-arrival[6] = 0.020
Inter-arrival[7] = 0.142
Inter-arrival[8] = 0.172
Inter-arrival[9] = 0.214
Inter-arrival[10] = 0.161
Inter-arrival[11] = 0.011
Inter-arrival[12] = 0.038
Inter-arrival[13] = 0.312
Inter-arrival[14] = 0.138
Inter-arrival[15] = 0.031
Inter-arrival[16] = 0.207
Inter-arrival[17] = 0.023
Inter-arrival[18] = 0.189
Inter-arrival[19] = 0.101
Inter-arrival[20] = 0.122
Inter-arrival[21] = 0.347
Inter-arrival[22] = 0.097
Inter-arrival[23] = 0.009
Inter-arrival[24] = 0.111
Inter-arrival[25] = 0.350
Inter-arrival[26] = 0.031
Inter-arrival[27] = 0.010
Inter-arrival[28] = 0.250
Inter-arrival[29] = 0.023
Inter-arrival[30] = 0.045
Inter-arrival[31] = 0.013
Inter-arrival[32] = 0.025
Inter-arrival[33] = 0.043
Inter-arrival[34] = 0.519
Inter-arrival[35] = 0.136
Inter-arrival[36] = 0.011
Inter-arrival[37] = 0.069
Inter-arrival[38] = 0.055
Inter-arrival[39] = 0.069
Inter-arrival[40] = 0.015
Inter-arrival[41] = 0.015
Inter-arrival[42] = 0.188
Inter-arrival[43] = 0.035
Inter-arrival[44] = 0.009
Inter-arrival[45] = 0.006
Inter-arrival[46] = 0.003
Inter-arrival[47] = 0.012
Inter-arrival[48] = 0.017
Inter-arrival[49] = 0.193
Inter-arrival[50] = 0.033
Inter-arrival[51] = 0.021
Inter-arrival[52] = 0.398
Inter-arrival[53] = 0.141
Inter-arrival[54] = 0.237
Inter-arrival[55] = 0.191
Inter-arrival[56] = 0.015
Inter-arrival[57] = 0.216
Inter-arrival[58] = 0.068
Inter-arrival[59] = 0.227
Inter-arrival[60] = 0.036
Inter-arrival[61] = 0.082
Inter-arrival[62] = 0.054
Inter-arrival[63] = 0.075
Inter-arrival[64] = 0.013
Inter-arrival[65] = 0.186
Inter-arrival[66] = 0.043
Inter-arrival[67] = 0.274
Inter-arrival[68] = 0.000
Inter-arrival[69] = 0.021
Inter-arrival[70] = 0.156
Inter-arrival[71] = 0.050
Inter-arrival[72] = 0.013
Inter-arrival[73] = 0.367
Inter-arrival[74] = 0.030
Inter-arrival[75] = 0.015
Inter-arrival[76] = 0.179
Inter-arrival[77] = 0.041
Inter-arrival[78] = 0.145
Inter-arrival[79] = 0.081
Inter-arrival[80] = 0.229
Inter-arrival[81] = 0.079
Inter-arrival[82] = 0.039
Inter-arrival[83] = 0.116
Inter-arrival[84] = 0.006
Inter-arrival[85] = 0.132
Inter-arrival[86] = 0.001
Inter-arrival[87] = 0.012
Inter-arrival[88] = 0.029
Inter-arrival[89] = 0.104
Inter-arrival[90] = 0.070
Inter-arrival[91] = 0.205
Inter-arrival[92] = 0.011
Inter-arrival[93] = 0.023
Inter-arrival[94] = 0.062
Inter-arrival[95] = 0.243
95
The expected value of interaarival time is 0.10
Inter-arrival times written to inter_arrival.txt
\end{verbatim}

\begin{figure}[H]
    \centering
    \includegraphics[width=0.5\textwidth]{Results/InterarrivalPicture.png}
    \caption{Inter arrival Time Modelling}
    \label{fig:placeholder}
\end{figure}
\noindent
\textbf{DISCUSSION:}\\
\noindent
The experiment helped in demonstrating that a Poisson arrival process produces exponentially distributed inter-arrival times. The simulated mean inter-arrival time closely matched the calculated theoretical value of $1/\lambda$, therefore validating the model. Due to randomness, some deviations were observed. The histogram of inter-arrival times showed the expected exponential distribution which is right skewed. \\

\noindent
\textbf{ACKNOWLEDGEMENT: }\\
I would like to share my thanks to Prof. Amitalok J Budkuley for explaining the context of the algorithms. I would like to extend my thanks to TAs Anuj Meena and Prince for questioning the fundamentals of the experiment and helping throughout along with the problem description. I would like to thank my colleagues Mayank Singh, Saptarshi Bhattacharjee and Soumava Biswas for collaborating in lab throughout the experiment. \\

\noindent
\printbibliography

\noindent
\textbf{APPENDIX (CODE): }\\
\appendix
\noindent
\textbf{Djikstra's Algorithm in C}

\begin{lstlisting}
#include <stdio.h>
#include <stdlib.h>
#include <math.h>
#include <time.h>

#define MAX_EVENTS 1000
float inter_arrival[MAX_EVENTS];
int ia_idx = 0;
// generating exponential RV
float exponential(float rate) {
    float u = (float)rand() / RAND_MAX;
    return -log(1.0 - u) / rate;
}

float expectation(float arr[], int n){
    
    float sum_ia = 0.0;
    for(int i = 0; i < n; i++){
        sum_ia += arr[i];
    }
    return sum_ia / n; 

}
int main() {

    srand(time(NULL));

    // parameters
    float lambda = 10.0;     // arrival rate
    float mu = 13.0;         // service rate
    int K = 5;               // finite buffer size
    float SIM_TIME = 10.0;  // total simulation time

    // variables
    int buffer = 0;
    int arrivals = 0;
    int dropped = 0;
    int serviced = 0;

    // event times
    float t = 0.0;
    float next_arrival = exponential(lambda);
    float next_service = INFINITY;
    float arrival_time[MAX_EVENTS];
    //float inter_arrival[1000];
    int arr_idx = 0;
    //int ia_idx = 0;

    // loop being less than the total simulation timew
    while (t < SIM_TIME) {

        if (next_arrival < next_service) {
            // arrival process
            t = next_arrival;
            arrival_time[arr_idx++] = t;
            arrivals++;

            if (buffer < K) {
                buffer++;

                //If server is idle, schedule service 
                if (buffer == 1) {
                    next_service = t + exponential(mu);
                }
            } else {
                dropped++;
            }

            next_arrival = t + exponential(lambda);

        } else {
            // service
            t = next_service;

            if (buffer > 0) {
                buffer--;
                serviced++;
            }

            if (buffer > 0) {
                next_service = t + exponential(mu);
            } else {
                next_service = INFINITY;
            }
        }
        //printf("Event Time next arrival : %0.2f\n", t);
        //printf("----------------------------------\n");
    }

    

    // result
    printf("Simulation Time           : %.2f\n", SIM_TIME);
    printf("Arrival rate (lambda)     : %.2f\n", lambda);
    printf("Service rate (mu)         : %.2f\n", mu);
    printf("Buffer size (K)           : %d\n", K);
    printf("----------------------------------\n");
    printf("Total Arrivals            : %d\n", arrivals);
    printf("Total Serviced            : %d\n", serviced);
    printf("Total Dropped             : %d\n", dropped);
    printf("Packet Drop Probability   : %f\n", (float)dropped / arrivals);
    printf("----------------------------------\n");
    
    
    printf("Storing Interaarival Times\n");
    for (int i = 1; i < arr_idx; i++) 
    {   
        inter_arrival[ia_idx] = arrival_time[i] - arrival_time[i-1];
        printf("Inter-arrival[%d] = %.3f\n", i, inter_arrival[ia_idx]);
        ia_idx++;
    }
    printf("%d\n", ia_idx);
    float expected_value = expectation(inter_arrival, ia_idx);
    printf("The expected value of interaarival time is %0.2f\n", expected_value);

    FILE *fp = fopen("inter_arrival.txt", "w");

    if (fp == NULL) {
        printf("Error opening file!\n");
        return 1;
    }

    for (int i = 0; i < ia_idx; i++) {
        fprintf(fp, "%.6f\n", inter_arrival[i]);
    }

    fclose(fp);

    printf("Inter-arrival times written to inter_arrival.txt\n");

    return 0;
}

\end{lstlisting}

\begin{lstlisting}[language=Python, caption=Python code for plotting the interarrival time]
import numpy as np
import matplotlib.pyplot as plt

data = np.loadtxt("inter_arrival.txt")

plt.hist(data, bins=50, density=True)
t = np.linspace(0, max(data), 100)
plt.plot(t, 10*np.exp(-10*t), 'r')
plt.show()
\end{lstlisting}
\end{document}
